\section{v1.2.0: readability and target trends}
The new behavior is explicit. Legacy legend layout and center-point text-color
estimation remain defaults. Data in this chapter are illustrative.
\subsection{Measured legend rows and columns}
\key{legend layout=measured} measures width, height and depth, using as many rows
or columns as needed. The whole legend can extend beyond the ramp length.
\key{legend duplicates=error} checks builtin numeric labels and reports both raw
ticks if formatting collapses them. Increase precision or reduce ticks; no values
are changed automatically. Custom formatters and siunitx are opaque and rejected
by this duplicate check. Measured layout accepts them and calls each tick formatter
once, reusing its box without extra Figure tags.
\begin{demo}{Measured numeric legends}
\gradheatlegend[min=-2,max=8,heat palette=purple-teal,
 heat center=1,legend ticks={-2,0,1,4,8},precision=1,
 width=60mm,height=3mm,unit={ pt},font=\large,
 legend layout=measured,legend duplicates=error,
 legend missing=false]
\qquad
\gradheatlegend[min=0,max=1,precision=1,
 legend ticks={0,.2,.4,.6,.8,1},font=\large,
 width=22mm,height=3mm,legend orientation=vertical,
 legend layout=measured,legend missing=false]
\end{demo}
Shared definitions capture both keys; \key{\gradlegenduse} accepts layout and
duplicate-check overrides. Defaults are \key{legacy} and \key{off}.
\clearpage
\subsection{Contrast over the label area}
\key{text contrast=area} applies to internal main labels with \key{text color=auto}.
It bounds the worst contrast over resolved fill, unfilled track and reference-band
intersections. Each gradient interval uses 32 channel-bounded subdivisions of
continuous sRGB luminance. It covers negative gradients and large labels;
\key{contrast minimum=4.5} remains explicit and is not lowered for large fonts.
\key{contrast failure=warn} reports an insufficient bound; \key{background} draws
an opaque white panel and black text. Explicit text colors bypass evaluation.
\begin{demo}{Readable internal labels}
\begingroup
\gradbarssetup{style=readable,width=62mm,
 min=-100,max=100,precision=0}
\begin{tabular}{lc}
Positive & \gradbar{80}\\
Negative & \gradbar{-65}\\
Zero & \gradbar{0}\\
Gradient & \gradbar[min=0,fill=gradient,
 left color=black,right color=white,
 text={Gradient label with panel}]{100}\\
\end{tabular}
\endgroup
\end{demo}
Graphic \key{style=readable} selects a rectangular 3ex accessible solid track,
small centered internal labels and area contrast with panel fallback. Later keys
override it. \key{contrast canvas=white} declares the background beyond the track;
set it to the actual cell/page color. The background is not inferred.
Rounded corners, patterns, compound shapes and overlaid reference/error strokes
are rejected by inside area evaluation. Inherited transparency, blending, custom
TikZ drawing, colored glyphs and font effects are outside the model. Heat and stack
segment labels keep their old algorithms. An existing opaque print-label panel
is assessed as white and covers underlying patterns, corners and reference strokes. This is not whole-document PDF/UA
certification. Use outside labels or explicit colors for unsupported drawings.
\clearpage
\subsection{Exactly one new combination: spark line plus target}
\key{\gradsparktarget} takes a sequence and a separate numeric target. It requires
\key{spark range=fixed} or a named scale and \key{spark type=line}. The target must
lie within the range, is not an observation and never participates in statistics.
NA preserves its position and breaks the line. Constant/single-point series use
the explicit scale; all-missing series retain the target and describe missing data.
\begin{demo}{Progress against an explicit target}
\begingroup
\gradscale{progress-target}{min=0,max=100}
\gradbarssetup{width=45mm,height=4ex,
 precision=0,unit={\%},label width=10mm}
\begin{tabular}{lc}
Team & Progress; dashed target: 70\%\\
\hline
Atlas & \gradsparktarget[scale group=progress-target]
 {20,35,NA,60,75}{70}\\
Birch & \gradsparktarget[scale group=progress-target]
 {50,50,50,50,50}{70}\\
Cedar & \gradsparktarget[scale group=progress-target]
 {NA,NA,NA,NA,NA}{70}\\
\end{tabular}
\endgroup
\end{demo}
The label uses the final position, even if missing. \key{target color} controls
the dashed horizontal line; the scalar bar \key{target} option stays rejected.
Automatic alternatives distinguish range, observations and reference, producing
one Figure. Existing \key{\gradspark} errors remain unchanged. Direct ordinary
\key{tabular} cells are supported; no new CSV compound record, native column type
or tabularray helper is added. Scalar backends need explicit TikZ fallback for
\key{spark-target} and \key{bar-area}.
\clearpage
\subsection{A readable native-table preset}
Table \key{style=readable} adds spacing, small type, accessible colors and a 22mm
outside-label slot. It preserves column semantics; internal panels require explicit
\key{options} selecting the graphic style.
\begin{demo}{Readable completion table}
\begin{gradtable}[style=readable]
\gradcolumn[width=24mm]{task}{Workstream}
\gradcolumn[type=bar,options={max=100,
 precision=0,unit={\%}}]{progress}{Completion}
\gradrow{Review,80}
\gradrow{Testing,35}
\gradrow{Pending,NA}
\end{gradtable}
\end{demo}

\par\smallskip\noindent Standalone source: \texttt{docs/examples/readability-target.tex}.
